\begin{lemma}[Scheide replacement tree is globally maximum (9)]
For each defective color $\delta$, the switched coloring
$\varphi_\delta$ is proper and leaves exactly $e_0$ uncolored, and the
replacement tree $T_\delta$ in the chosen replacement data forms a globally
maximum triple $(T_\delta,e_0,\varphi_\delta)$.
\end{lemma}
\begin{proof}
Put $r=\gammaIndex(\delta)$ and
$b=p-1+2r$, the zero-based position of the first edge in the block of
color $\gamma(r)$.  Write $n$ for the length of $T$.  The ordered-data
clauses of \noderef{ScheideReplacementData} say that $T_\delta$ also has
length $n$: positions below $b$ are retained, positions $j$ with
$b\leq j$ and $j+2<n$ contain the old entry at position $j+2$, and the
last two edge positions $n-2,n-1$ contain $f^1_\delta,f^2_\delta$.
The corresponding retained-vertex clauses keep every vertex position at
most $b$ and shift every old vertex position strictly above the deleted
pair down by two.  Here these positions are meaningful because
$r<s$, $p\geq2$, and
$n+1=p+2s$, all recorded by the inherited structural data in
\noderef{ScheideDefectiveSetup}, \noderef{ScheideReducedConfiguration},
and \noderef{ScheideStructuralConfiguration}.

First, \noderef{ScheideDefectiveSwitchEffect} gives, verbatim, that
$\varphi_\delta$ is proper and that exactly $e_0$ is uncolored.  It also
gives the three switch facts used below: every tree-incident edge other
than $f^0_\delta$ keeps its color; every tree vertex other than
$v^0_\delta$ keeps its missing-color set; and
\[
 \overline{\varphi_\delta}(v^0_\delta)
 =\bigl(\overline\varphi(v^0_\delta)
        \setminus\{\gamma(r)\}\bigr)\cup\{\delta\}.
\tag{*}
\]

We next verify \noderef{TashkinovTree} for $T_\delta$ one edge position
at a time.  Fix a noninitial edge position $j$, so $1\leq j<n$.
There are three exhaustive cases.

Suppose first that $j<b$.  The retained-entry equalities from
\noderef{ScheideReplacementData} identify the new edge at $j$ with the
old edge at $j$.  The globally maximum field inherited from
\noderef{ScheideStructuralConfiguration}, interpreted with
\noderef{GlobalMaximumTriple}, says that the old tree is Tashkinov.
Consequently \noderef{TashkinovTree} supplies an index $a\leq j$ and a
color $c$ such that the old edge at $j$ has color $c$ and $c$ is missing
at the old vertex at $a$.  Since $a\leq j<b$, the retained-vertex equality
identifies that old witness with the new vertex at the same index $a$.

The old edge at $j$ is incident with $V(T)$ and is not
$f^0_\delta$.  Indeed, the attachment clause of \noderef{TreeSequence}
makes both ends of a tree edge vertices listed in $T$, hence members of
the image defining \noderef{TreeSequenceVertexSet}; but
\noderef{ScheideDefectiveSetup} gives
$f^0_\delta=v^0_\delta u^0_\delta$ with
$u^0_\delta\notin V(T)$.  Thus the switch-effect clause preserves the
edge color $c$.

If the witness vertex at $a$ is not $v^0_\delta$, its missing set is also
unchanged.  If it is $v^0_\delta$, we must check that $(*)$ does not erase
$c$.  When $j<p-1$, the edge is a nonfirst base edge because $j\geq1$;
the nonfirst-base-edge avoidance clause of
\noderef{ScheideStructuralConfiguration} gives
$c\ne\gamma(r)$.  When $p-1\leq j<b$, division of
$j-(p-1)$ by two writes the entry uniquely as the first or second edge of
a block with index $h<r$.  Its block-color clause gives $c=\gamma(h)$,
and injectivity of the indexed block colors gives
$c\ne\gamma(r)$.  In either subcase $(*)$ preserves $c$ as missing at the
same vertex.  Hence $a\leq j$ is a valid earlier witness for the retained
edge.

Suppose next that $b\leq j$ and $j+2<n$.  The retained-after equality in
\noderef{ScheideReplacementData} identifies the new edge at $j$ with the
old edge at $j+2$.  Since $j+2\geq b+2$, the exact layout in
\noderef{ScheideStructuralConfiguration} writes that old entry as one of
the two edges of a unique block $h>r$.  The block-color clause colors it
$\gamma(h)$, and the block-color-missing-on-base clause supplies a base
index $a<p$ at whose vertex $\gamma(h)$ is missing.  The inequalities
\[
 a\leq p-1\leq b\leq j
\]
show both that the replacement data retain this base vertex at index $a$
and that it is earlier than the edge at zero-based position $j$ in the
sense required by \noderef{TashkinovTree}.  The old block edge is again a
tree edge distinct from $f^0_\delta$, by the same endpoint argument, so
its color is fixed by \noderef{ScheideDefectiveSwitchEffect}.  If the
chosen base vertex is not $v^0_\delta$, its missing set is fixed as well;
if it is $v^0_\delta$, then $h>r$ and block-color injectivity give
$\gamma(h)\ne\gamma(r)$, so $(*)$ still leaves $\gamma(h)$ missing.
Thus the shifted edge also has the required retained base witness.

It remains that $j+2\geq n$.  Since $j<n$, this says precisely that
$j=n-2$ or $j=n-1$, the two appended positions.  The exit clause in
\noderef{ScheideDefectiveSetup} places $v^0_\delta$ among the Scheide exit
vertices, and \noderef{ScheideStructuralExitsInBase} therefore gives an
index $a<p$ with $v^0_\delta=y_a$.  The base-position clause of
\noderef{ScheideStructuralConfiguration}, followed by the
retained-vertices-before clause of \noderef{ScheideReplacementData},
identifies the replacement vertex at index $a$ with $v^0_\delta$; the
required retention inequality is $a\leq p-1\leq b$.  Moreover $r<s$
implies $s\geq1$, so $n+1=p+2s$ implies
\[
 a\leq p-1\leq n-2<n-1.
\]
Thus this witness precedes both appended edge positions.  Formula $(*)$
says that $\delta$ is missing there.

For the first appended edge, the endpoint equality
$f^1_\delta=v^1_\delta u^1_\delta$ in
\noderef{ScheideReplacementData} makes it incident with
$u^1_\delta$.  Hence \noderef{ScheideReplacementPairFixed} and the
recorded equality $\varphi(f^1_\delta)=\delta$ give
$\varphi_\delta(f^1_\delta)=\delta$.  The same argument using
$f^2_\delta=v^2_\delta u^2_\delta$ and
$\varphi(f^2_\delta)=\delta$ gives
$\varphi_\delta(f^2_\delta)=\delta$.  Therefore the retained occurrence
of $v^0_\delta$ at index $a$ witnesses the Tashkinov condition for each
of the last two edges.  The three cases prove that $T_\delta$ is a
Tashkinov tree under $\varphi_\delta$.

We now compute its order without using the replacement vertex-set
equation.  By \noderef{TreeSequenceVertexSet}, the vertex set of any
tree-sequence $S$ is the image of all indices
$0,\ldots,S.\mathrm{length}$ under its vertex map.  That map is injective
by \noderef{TreeSequence}, so the image has cardinality
$S.\mathrm{length}+1$.  Applying this once to $T_\delta$ and once to $T$,
and then using the replacement-length equality from
\noderef{ScheideReplacementData}, gives
\[
 |V(T_\delta)|=n+1=|V(T)|.
\]
The cardinality clause of the old globally maximum triple in
\noderef{GlobalMaximumTriple} says $|V(T)|=t(G)$; hence
$|V(T_\delta)|=t(G)$.

Finally unfold \noderef{KTriple}.  From the $K$-triple clause of the old
globally maximum field we retain, unchanged, that $G$ is critical, that
$\chi'(G)=k+1$, and that $\Delta(G)+1\leq k$.  The first paragraph supplied
properness of $\varphi_\delta$ and the assertion that exactly $e_0$ is
uncolored.  These five clauses assemble the new $K$-triple.  Combining it
with the proved Tashkinov-tree clause and the proved equality
$|V(T_\delta)|=t(G)$ gives precisely the three defining clauses of
\noderef{GlobalMaximumTriple}.  This proves the claimed conjunction.
\end{proof}
